% Part: first-order-logic
% Chapter: natural-deduction
% Section: derivations

\documentclass[../../../include/open-logic-section]{subfiles}

\begin{document}

\iftag{FOL}
      {\olfileid[fr]{fol}{ntd}{der}}
      {\olfileid[fr]{pl}{ntd}{der}}

\olsection{\usetoken{P}{derivation}}

\begin{explain}
Nous avons défini les hypothèses et donné les règles d'inférence.
Les !!{derivation}s en déduction naturelle sont engendrées
inductivement à partir de ces éléments : chaque !!{derivation}
est soit une hypothèse isolée, soit constituée d'une, de deux ou de
trois !!{derivation}s suivies d'une inférence correcte.
\end{explain}

\begin{defn}[!!^{derivation}]
\Article{derivation} \emph{!!{derivation}} d'un !!{sentence}~$!A$
à partir d'un ensemble d'hypothèses~$\Gamma$ est un arbre fini
d'!!{sentence}s qui satisfait les conditions suivantes :
\begin{enumerate}
\item Les !!{sentence}s situés tout en haut de l'arbre appartiennent
  à $\Gamma$ ou sont déchargés par une inférence de l'arbre.
\item L'!!{sentence} situé tout en bas de l'arbre est~$!A$.
\item Chaque !!{sentence} de l'arbre, à l'exception de
  l'!!{sentence}~$!A$ situé tout en bas, est une prémisse d'une
  application correcte d'une règle d'inférence dont la conclusion
  se trouve immédiatement en dessous de cet !!{sentence}.
\end{enumerate}
On dit alors que $!A$ est la \emph{conclusion} de la !!{derivation}
et que $\Gamma$ est son contexte d'hypothèses !!{undischarged}s.
Les hypothèses effectivement non déchargées doivent appartenir à
$\Gamma$, sans nécessairement épuiser cet ensemble.

S'il existe une !!{derivation} de $!A$ à partir de~$\Gamma$,
on dit que $!A$ est \emph{!!{derivable}} à partir de~$\Gamma$,
ce que l'on note $\Gamma \Proves !A$. S'il existe une
!!{derivation} de~$!A$ dans laquelle toutes les hypothèses sont
!!{discharged}s, on écrit~$\Proves !A$.
\end{defn}

\begin{ex}
Toute hypothèse isolée est une !!{derivation}. Ainsi, $!A$ seul
constitue une !!{derivation}, et il en va de même pour $!B$ seul.
On peut en obtenir une nouvelle !!{derivation} en appliquant,
par exemple, la règle $\Intro{\land}$ :
\begin{prooftree}
\AxiomC{$!A$}
\AxiomC{$!B$}
\RightLabel{\Intro{\land}}
\BinaryInfC{$!A \land !B$}
\end{prooftree}
Ces règles sont générales : on peut y remplacer $!A$ et~$!B$
par des !!{sentence}s quelconques, par exemple $!C$ et~$!D$.
La conclusion est alors $!C \land !D$, et l'arbre
\begin{prooftree}
\AxiomC{$!C$}
\AxiomC{$!D$}
\RightLabel{\Intro{\land}}
\BinaryInfC{$!C \land !D$}
\end{prooftree}
est une !!{derivation} correcte. On peut aussi intervertir les
hypothèses, de sorte que $!D$ joue le rôle de~$!A$, et $!C$ celui
de~$!B$. L'arbre
\begin{prooftree}
\AxiomC{$!D$}
\AxiomC{$!C$}
\RightLabel{\Intro{\land}}
\BinaryInfC{$!D \land !C$}
\end{prooftree}
est donc lui aussi une !!{derivation} correcte.

On peut maintenant appliquer une autre règle, par exemple
$\Intro{\lif}$. Elle permet de conclure une implication et de
!!{discharge} toute hypothèse identique à l'antécédent de cette
implication. Les deux !!{derivation}s suivantes sont donc correctes :
\begin{prooftree}
\AxiomC{$\Discharge{!C}{1}$}
\AxiomC{$!D$}
\RightLabel{\Intro{\land}}
\BinaryInfC{$!C \land !D$}
\DischargeRule{\Intro{\lif}}{1}
\UnaryInfC{$!C \lif (!C \land !D)$}
\DisplayProof\bottomAlignProof
\AxiomC{$!C$}
\AxiomC{$\Discharge{!D}{1}$}
\RightLabel{\Intro{\land}}
\BinaryInfC{$!C \land !D$}
\DischargeRule{\Intro{\lif}}{1}
\UnaryInfC{$!D \lif (!C \land !D)$}
\end{prooftree}
Elles montrent respectivement que $!D \Proves !C \lif (!C \land !D)$
et que $!C \Proves !D \lif (!C \land !D)$.

Rappelons que la décharge d'hypothèses est permise, mais non
obligatoire. En particulier, une règle reste applicable même si
les hypothèses qu'elle permet de décharger ne figurent pas dans
la !!{derivation}. L'inférence suivante est ainsi correcte,
bien qu'il n'y ait aucune hypothèse~$!A$ à !!{discharge} :
\begin{prooftree}
  \AxiomC{$!B$}
  \DischargeRule{\Intro{\lif}}{1}
  \UnaryInfC{$!A \lif !B$}
\end{prooftree}
\end{ex}

\begin{editorial}
La source appelle $\Gamma$ les hypothèses non déchargées. La
précision sur le contexte explicite sa propre définition, qui exige
seulement que les hypothèses non déchargées appartiennent à
$\Gamma$, sans imposer que tous ses éléments soient utilisés.
\end{editorial}

\end{document}
